% from ega2-1-fr.tex
\oldpage[II]{93}

soient affines et recouvrent $Y$
(\hyperref[II.4.5.2-fr]{4.5.2, a')})~;
remplaçant les $s_i$ par des puissances convenables, on peut en outre
supposer tous les $k_i$ égaux à un même entier $k$. Soient $s'_i$ les
sections de $f^*(\mathcal{K}^{\otimes k})$ au-dessus de $X$ correspondant
canoniquement aux $s_i$, de sorte que les
$X_{s'_i}=f^{-1}(Y_{s_i})$ (4.6.1.13, (i))
recouvrent $X$. Comme $\mathcal{L}|X_{s'_i}$ est ample
((\hyperref[II.4.6.4-fr]{4.6.4}) et
(\hyperref[II.4.6.6-fr]{4.6.6})), il existe pour chaque $i$ des sections en
nombre fini $t_{ij}\in\Gamma(X,\mathcal{L}^{\otimes n_{ij}})$ tels que les
$X_{t_{ij}}$ soient affines, contenus dans $X_{s'_i}$ et recouvrent
$X_{s'_i}$ (\hyperref[II.4.5.2-fr]{4.5.2, a')})~; on peut d'ailleurs
supposer tous les $n_{ij}$ égaux à un même nombre $n$. Cela étant, $X$ est
séparé et quasi-compact, donc il existe un entier $m>0$, et pour tout
$(i,j)$ une section
\[
  u_{ij}\in\Gamma(X,\mathcal{L}^{\otimes n}\otimes
  f^*(\mathcal{K}^{\otimes mk}))
\]
tels que $t_{ij}\otimes s_i'^{\otimes m}$ soit la restriction à $X_{s'_i}$
de $u_{ij}$ (\hyperref[I.9.3.1-fr]{I, 9.3.1})~; en outre on a
$X_{u_{ij}}=X_{t_{ij}}$, donc les $X_{u_{ij}}$ sont affines et recouvrent
$X$. On peut d'ailleurs supposer que $m$ est de la forme $nr$~; si on pose
$n_0=rk$, on voit donc (\hyperref[II.4.5.2-fr]{4.5.2, a')}) que
$\mathcal{L}\otimes_{\mathcal{O}_X}f^*(\mathcal{K}^{\otimes n_0})$ est
ample. En outre,
il existe $h_0>0$ tel que $\mathcal{K}^{\otimes h}$ soit engendré par ses
sections au-dessus de $Y$ dès que $h\geq h_0$
(\hyperref[II.4.5.5-fr]{4.5.5})~; a fortiori,
$f^*(\mathcal{K}^{\otimes h})$ est engendré par ses sections au-dessus de
$X$ pour $h\geq h_0$, par définition des images réciproques
((\hyperref[0.3.7.1-fr]{0, 3.7.1}) et
(\hyperref[II.4.4.1-fr]{4.4.1})). On en conclut que
$\mathcal{L}\otimes f^*(\mathcal{K}^{\otimes n_0+h})$ est ample dès que
$h\geq h_0$ (\hyperref[II.4.5.6-fr]{4.5.6}), ce qui achève la
démonstration.

\begin{remark}[4.6.14]
\label{II.4.6.14-fr}
Sous les conditions de (ii), on se gardera de croire que
$\mathcal{L}\otimes f^*(\mathcal{K})$ soit ample pour $g\circ f$~; en effet,
comme $\mathcal{L}\otimes f^*(\mathcal{K}^{-1})$ est aussi ample pour $f$
(\hyperref[II.4.6.5-fr]{4.6.5}), on en conclurait que $\mathcal{L}$ serait
ample pour $g\circ f$~; prenant en particulier pour $g$ le morphisme
identique, tout $\mathcal{O}_X$-Module inversible serait ample pour $f$, ce
qui n'est pas le cas en général (voir (\hyperref[II.5.1.6-fr]{5.1.6}),
(\hyperref[II.5.3.4-fr]{5.3.4, (i)}) et
(\hyperref[II.5.3.1-fr]{5.3.1})).
\end{remark}

\begin{proposition}[4.6.15]
\label{II.4.6.15-fr}
Soient $f:X\to Y$ un morphisme quasi-compact, $\mathcal{J}$ un faisceau
quasi-cohérent d'idéaux localement nilpotent de $\mathcal{O}_X$, $Z$ le
sous-préschéma fermé de $X$ défini par $\mathcal{J}$, $j:Z\to X$
l'injection canonique. Pour qu'un $\mathcal{O}_X$-Module inversible
$\mathcal{L}$ soit ample pour $f$, il faut et il suffit que
$j^*(\mathcal{L})$ soit ample pour $f\circ j$.
\end{proposition}

En effet, la question étant locale sur $Y$
(\hyperref[II.4.6.4-fr]{4.6.4}), on peut supposer $Y$ affine~; $X$ étant
alors quasi-compact, on peut supposer $\mathcal{J}$ nilpotent. Compte tenu
de (\hyperref[II.4.6.6-fr]{4.6.6}), la proposition n'est autre alors que
(\hyperref[II.4.5.13-fr]{4.5.13}).

\begin{corollary}[4.6.16]
\label{II.4.6.16-fr}
Soient $X$ un préschéma localement noethérien, $f:X\to Y$ un morphisme
quasi-compact. Pour qu'un $\mathcal{O}_X$-Module inversible $\mathcal{L}$
soit ample pour $f$, il faut et il suffit que son image réciproque
$\mathcal{L}'$ par l'injection canonique $X_{\mathrm{red}}\to X$ soit ample
pour $f_{\mathrm{red}}$.
\end{corollary}

On a déjà vu que la condition est nécessaire
(\hyperref[II.4.6.13-fr]{4.6.13, (vi)})~; inversement, si elle est remplie,
on peut se borner, pour prouver que $\mathcal{L}$ est ample pour $f$, au
cas où $Y$ est affine (\hyperref[II.4.6.4-fr]{4.6.4})~; alors
$Y_{\mathrm{red}}$ est aussi affine, donc $\mathcal{L}'$ est ample
(\hyperref[II.4.6.6-fr]{4.6.6}), et il en est de même de $\mathcal{L}$ en
vertu de (\hyperref[II.4.5.13-fr]{4.5.13}), puisque alors $X$ est
noethérien et $X_{\mathrm{red}}$ un sous-préschéma fermé de $X$ défini par
un faisceau quasi-cohérent nilpotent d'idéaux
(\hyperref[I.6.1.6-fr]{I, 6.1.6}).

\begin{proposition}[4.6.17]
\label{II.4.6.17-fr}
Les notations et hypothèses étant celles de
(\hyperref[II.4.4.11-fr]{4.4.11}), pour que $\mathcal{L}''$ soit ample
relativement à $f''$, il faut et il suffit que $\mathcal{L}$ soit ample
relativement à $f$ et $\mathcal{L}'$ ample relativement à $f'$.
\end{proposition}

